\section{El CEO como puente: investigación, comercialización y expectativas de la organización}

Antes y después de la salida a bolsa, el papel del CEO cambia. Zhang Peng (张鹏) dice que ahora dedica más energía a la operación diaria de la empresa, en especial al área de atención al cliente, la orientación al mercado y la comercialización. También menciona que el CEO muchas veces es un puente, quien monta el escenario: hace que los departamentos de investigación y de tecnología y los equipos de comercialización puedan comunicarse, y que todos desplieguen su imaginación, su capacidad y su capacidad de ejecución.

\subsection{La investigación y la comercialización no pueden avanzar desconectadas}

En una empresa de modelos es fácil que haya dos piezas desconectadas: el equipo de investigación se ocupa de las capacidades del modelo y el equipo comercial, de las entregas al cliente. Zhang Peng subraya que internamente hay que alinear de forma continua los cambios del mercado, la investigación y el desarrollo tecnológicos y las necesidades de comercialización. De lo contrario, la investigación puede producir algo vistoso pero difícil de llevar a la práctica, y el área comercial puede sacrificar la dirección tecnológica de largo plazo a cambio de ingresos de corto plazo.

\begin{importantbox}{La doble tarea del CEO de una empresa de modelos}
Por un lado, debe proteger la investigación de largo plazo y las capacidades del modelo; por otro, debe asegurar el valor para el cliente, los ingresos, el cumplimiento normativo y la eficiencia de la organización. Lo central del CEO no es decidir por todos, sino lograr que estos sistemas puedan escucharse entre sí.
\end{importantbox}

\subsection{Algunos momentos conmovedores}

El papel de puente del CEO a veces resulta abstracto, por eso esta sección lo aterriza con algunos momentos concretos. Muestran que una empresa de modelos no se demuestra solo en artículos y presentaciones de lanzamiento.

Al final de la entrevista, Zhang Peng recuerda algunos momentos. Uno es haber pasado casi medio año en Shenzhen por un cliente grande y, al final, no volver a Pekín con las manos vacías, sino con contratos por decenas de millones; esto muestra que los clientes están dispuestos a pagar por la tecnología. Otro es haber repartido sobres rojos a los asistentes con un agente inteligente durante una presentación de lanzamiento: aunque el monto se escribió con un dígito equivocado, los sobres rojos sí se enviaron, y se comentó que era el primer sobre rojo que la IA le daba a la humanidad. Otro más es que, después de que GLM-4.5 se publicó como código abierto, desarrolladores y productos del extranjero empezaron a usarlo, e incluso aparecieron toda clase de reutilizaciones, destilaciones y reempaquetados.

Estos momentos muestran en conjunto una misma cosa: la satisfacción de una empresa de modelos no viene solo de los benchmark, sino también de los contratos con clientes, el uso real, el reconocimiento de los desarrolladores y la repercusión de su tecnología.

\begin{knowledgebox}{Del benchmark a la repercusión real}
Un benchmark demuestra capacidad; un contrato con un cliente demuestra valor; el uso del código abierto demuestra ecosistema; el evento de lanzamiento demuestra experiencia. Una empresa de modelos tiene que gestionar al mismo tiempo estos distintos tipos de retroalimentación.
\end{knowledgebox}

\subsection{Resumen de la sección}

El relato de Zhang Peng como CEO no busca el «estrellato»; se parece más a un papel de puente: conecta la investigación, la ingeniería, el área comercial, los clientes, la gobernanza y las expectativas del equipo. Esta es también la otra cara de la «sensación de cemento» de Zhipu (智谱).
